\section{Background and Related Work}
\label{sec:background}

\subsection{Live-cell and in-incubator imaging}
Live-cell imaging observes cultures over time under controlled temperature,
humidity and CO\textsubscript{2}. Commercial in-incubator systems (e.g.\ IncuCyte-class
instruments) automate multi-well time-lapse acquisition but are closed and
costly. This thesis targets the specific case of \emph{perfused bioreactor
chips}, where the sample is a sealed microfluidic device observed through a
small optical window rather than an open well plate.
\todo{cite the specific commercial baseline and the bioreactor-chip platform
used in the wet-lab experiments.}

\subsection{Open-hardware microscopy}
The OpenFlexure Microscope \cite{Collins2020,Sharkey2016} is a 3D-printed,
flexure-based motorised microscope driven by a Raspberry~Pi and a Pi camera; it
provides sub-micron mechanical stepping via a monolithic flexure stage and an
established software stack (OpenFlexure Server). Related open platforms include
UC2 \cite{Diederich2020} (modular cube-based optics) and the Squid/Octopi
family \cite{Li2020Squid} (higher-performance automated imaging). The mesoscope
reuses the OpenFlexure optics and Z-focus mechanics but departs from the XY
flexure stage: it removes lateral flexure motion and instead tilts the whole
body and translates it along an external linear rail
(Section~\ref{sec:mechanical}).

\subsection{Edge AI on neural-processing units}
On-device inference avoids transferring images off the instrument and enables
near-real-time analysis inside the incubator. The Raspberry~Pi AI~HAT+ carries a
Hailo-8 NPU (\SI{26}{TOPS}); models must be compiled to the Hailo Executable
Format (HEF) with the Hailo Dataflow Compiler and quantised to run on the NPU.
\todo{cite Hailo-8 architecture / dataflow compiler documentation.}

\subsection{Cell detection and segmentation}
Candidate model families for counting and delineating cells include
general-purpose detectors and instance segmenters (YOLOv8 and variants) and
cell-specific segmentation models such as Cellpose \cite{Stringer2021Cellpose}
and StarDist \cite{Schmidt2018StarDist}. The choice between bounding-box
detection and instance segmentation---and thus the CVAT label schema---depends
on the biological readout (confluence, count, morphology, migration), which is
an open decision fixed in milestone~M1
(Section~\ref{sec:requirements}). Only a foreign YOLOv8s-Pose benchmark
currently exists in the project; the task-specific model is built in
Section~\ref{sec:edgeai}.
